% Draft status: V0
% Evidence status: checked where possible; TODOs mark missing evidence.
\section{Introduction}

Object-recognition models are often compared with aggregate accuracy, but aggregate accuracy can hide uneven performance across regions. This is a concern for geographically diverse visual data because the same object category can appear in different environments, styles, materials, backgrounds, and socioeconomic contexts. Prior work on Dollar Street suggests that object-recognition performance may vary across countries and income levels \citep{devries2019object}. This motivates evaluating models not only by overall object accuracy but also by their best- and worst-region performance.

This paper studies regional bias in computer vision models on GeoDE and Dollar Street. GeoDE is a geographically diverse evaluation dataset for object recognition \citep{ramaswamy2023geode}. Dollar Street provides household-object imagery intended to represent geographic and socioeconomic visual diversity \citep{rojas2022dollarstreet}. These datasets make it possible to ask whether strong object recognition performance is accompanied by stable performance across regions.

The main research question is: how do pretrained and finetuned vision models differ in regional bias across GeoDE and Dollar Street? The paper compares nine model names: CLIP, ConvNeXt, DeiT, DINO, MAE, MaxViT, ResNet, SigLIP, and Swin. It evaluates two settings: pretrained ImageNet evaluation and finetuned evaluation. The analysis uses completed result files and derived tables; it does not train new models inside the paper workspace.

The central metrics are object accuracy, best-region accuracy, worst-region accuracy, and regional bias gap. The regional bias gap is the difference between best-region and worst-region accuracy. This metric is simple, but it helps prevent a purely aggregate reading of model performance. A model can have high object accuracy and still show a large regional gap. A model can also have a small regional gap because it performs poorly everywhere. For that reason, this paper interprets object accuracy and regional gap together.

The first draft results suggest three broad patterns. First, finetuning is associated with much higher object accuracy for every model-dataset pair in the derived table. Second, finetuning is also associated with larger regional bias gaps in every model-dataset pair. Third, the dataset matters: finetuned GeoDE results have higher mean object accuracy and lower mean regional gap than finetuned Dollar Street results. These are descriptive findings and should not be read as causal claims.

The contribution of this paper is a reproducible first analysis of regional object-recognition performance across two datasets, nine models, and two evaluation settings. It also provides a workflow for separating empirical claims from literature claims. Empirical claims come from normalized result tables and figures; literature claims come from verified citation entries. TODO: sharpen contributions after the related work and dataset audit are expanded.

